\section{Discussion}
\label{sec:discussion}

\textbf{Where the gain comes from.} The controlled ablation resolves the
design space of Section~\ref{sec:objectives}: training-side alternatives
(Focal Loss, many-hot supervision) leave the BCE control statistically
unchanged, whereas post-processing moves the operating point --- Gaussian
smoothing changes boundary decisions and temperature scaling improves
probability quality without changing the ranking. For frozen-feature pipelines this ordering matters
practically, because the post-processing search costs CPU milliseconds per
million frames while retraining alternatives cost GPU hours, and because the
search can be re-run per deployment domain without touching the
representation.

\textbf{What aggregate calibration hides.} After temperature scaling the
equal-width ECE falls to 0.002214, yet the class-balanced ECE remains
0.070405 --- roughly thirty times larger. The gap is a direct consequence of
boundary sparsity: non-boundary frames dominate every confidence bin, so a
near-perfect aggregate score is compatible with poorly calibrated
probabilities exactly on the rare positive class that the task cares about.
Reporting only aggregate ECE would therefore overstate calibration quality
for SBD and, we suspect, for other sparse temporal-event tasks. Class-wise
or event-conditional calibration (for example, fitting a separate temperature
on boundary-adjacent frames) is a natural next step, but it requires an
independent calibration split that the present protocol does not reserve.

\textbf{Sources of uncertainty.} The two uncertainty analyses answer
complementary questions. Paired video bootstrap intervals
(Table~\ref{tab:paired_delta}) quantify how much the measured deltas depend
on which videos were sampled, and remain the dominant term: their widths are
an order of magnitude larger than the seed-induced standard deviations of
Section~\ref{sec:seed_robustness}. Training-seed variability, the usual
objection to single-run reports, is bounded below 0.005 F1 across all three
datasets by the three-seed study. The remaining unquantified term is protocol
variability --- the frozen five-fold selection is itself a function of the
validation videos --- which could be probed by repeating the selection over
bootstrap-resampled validation sets.

\textbf{The ClipShots gap is a precision problem.} The transition-type
breakdown shows gradual-transition recall (0.8089) trailing cut recall
(0.8991), but the dominant loss is aggregate precision
(\PaperDeployClipPrecision): open-world content produces false-positive
peaks that survive smoothing. A frozen
representation cannot adapt to this distribution shift by construction, so
threshold or temperature adjustments only trade recall for precision along a
fixed curve. Closing the gap likely requires domain-adaptive components ---
either lightweight adapters on the frozen backbone or domain-specific
post-processing selection --- rather than further tuning of the present
search space.
